\documentclass[11pt, a4paper]{article}

% --- ESSENTIAL PACKAGES ---
\usepackage[T1]{fontenc}
\usepackage[utf8]{inputenc}
\usepackage{geometry}
\geometry{
    a4paper,
    left=2.5cm,
    right=2.5cm,
    top=3cm,
    bottom=3cm
}
\usepackage{xcolor}
\usepackage{hyperref}
\usepackage{enumitem}

% --- HYPERLINKS ---
\hypersetup{
    colorlinks=true,
    linkcolor=blue,
    urlcolor=blue,
    pdftitle={Project Application - Jakub Mitura},
}

% --- DOCUMENT ---
\begin{document}

\section*{Project Application - Jakub Mitura}

\section*{The Project}

\subsection*{Project details}

\noindent\textbf{Project title:} Generative AI for Nuclear Medicine Optimization (CausalPCa)

\vspace{0.5em}
\noindent\textbf{Project summary (abstract):} \\
This project introduces a new paradigm in oncology with an interactive, Generative AI (GenAI)-based agent for managing prostate cancer (PCa). We will build a causal "digital twin" to transparently model the disease trajectory by integrating complex multimodal data into a unified patient view. Technologically, we will fuse multimodal data, use GenAI for data augmentation, and create a medical knowledge base by disentangling disease signals from confounders. Clinically, our agent will improve diagnosis by simulating disease progression and enable personalized treatment selection by forecasting outcomes. The framework is designed as a trustworthy, explainable, and safe AI system compliant with the EU AI Act, establishing a new European standard for medical GenAI.

\vspace{0.5em}
\noindent\textbf{Keywords:} \\
Medical AI, Generative AI, Causal Inference, Prostate Cancer, Oncology, Digital Twin, Neural Jump ODE.

\vspace{0.5em}
\noindent\textbf{Instructions:} Not provided

\vspace{0.5em}
\noindent\textbf{Proposal for civilian purposes:} Yes.

\vspace{0.5em}
\noindent\textbf{Is any part of the project confidential?:} No.

\vspace{0.5em}
\noindent\textbf{Does your proposal involve handling of personal data?:} Yes.

\vspace{0.5em}
\noindent\textbf{Instructions:} Not provided

\subsection*{Submission details}

\noindent\textbf{Type of submission:} New Proposal

\subsection*{Industry, academia and public sector involvement}

\noindent\textbf{Industry involvement:} None.

\vspace{0.5em}
\noindent\textbf{Academic involvement:} Yes (Universitätsklinik für Radiologie und Nuklearmedizin Magdeburg).

\vspace{0.5em}
\noindent\textbf{Public sector involvement:} Yes (University Hospitals).

\section*{Principal Investigator}

\subsection*{Personal information}

\noindent\textbf{Gender:} Male \\
\textbf{Title:} Prof. Dr. \\
\textbf{First (given) name:} Michael \\
\textbf{Last (family) name:} Kreißl \\
\textbf{Initials:} MK \\
\textbf{Date of birth:} [Confidential] \\
\textbf{ID:} EHPC-AIF-2026SC01-062

\vspace{0.5em}
\noindent\textbf{1/7}

\vspace{0.5em}
\noindent\textbf{E-mail address:} [Email of PI] \\
\textbf{Secondary e-mail address:} [Secondary Email] \\
\textbf{Nationality:} German \\
\textbf{Phone number:} [Phone Number] \\
\textbf{Job title:} Head of Nuclear Medicine \\
\textbf{Employment contract valid for more than 3 months after end allocation:} Yes \\
\textbf{Website:} https://www.med.uni-magdeburg.de/kln/

\subsection*{Organization details}

\noindent\textbf{Instructions:} Not provided

\vspace{0.5em}
\noindent\textbf{Organization name:} Universitätsklinik für Radiologie und Nuklearmedizin Magdeburg (UKMD) \\
\textbf{Organization type:} Higher Education / Research Organisation \\
\textbf{Organization with research activity:} Yes \\
\textbf{Organization head office is located in Europe:} Yes \\
\textbf{Percentage of R\&D in Europe vs total R\&D:} 100\% \\
\textbf{Organization department:} Klinik für Radiologie und Nuklearmedizin \\
\textbf{Organization group:} Nuclear Medicine Research Group \\
\textbf{Organization address:} Leipziger Str. 44 \\
\textbf{Organization postal code:} 39120 \\
\textbf{Organization city:} Magdeburg \\
\textbf{Organization country:} Germany

\section*{Track record of the PI}

\noindent\textbf{Granted patents and other measures for the relevance of the work:} \\
Not provided

\vspace{0.5em}
\noindent\textbf{Prior allocation history in EuroHPC, PRACE, national calls, as well as international programs such as INCITE of the US DoE:} \\
Previous NHR Normal allocations for "CausalPCa" development (96,000 GPU hours).

\vspace{0.5em}
\noindent\textbf{Participation of team members in other European Commission (EC) actions, such as ERC or Marie SkłodowskaCurie EC grants, etc.:} \\
Not provided

\vspace{0.5em}
\noindent\textbf{ID:} EHPC-AIF-2026SC01-062

\vspace{0.5em}
\noindent\textbf{2/7}

\vspace{0.5em}
\noindent\textbf{Previous presentations at EuroHPC Summit, EuroHPC User day or PRACE days:} \\
Not provided

\section*{Contact Person and Team Members}

\subsection*{Contact person}

\noindent\textbf{First (given) name:} Jakub \\
\textbf{Last (family) name:} Mitura \\
\textbf{E-mail address:} [Email of Contact]

\vspace{0.5em}
\noindent\textbf{divider:} Not provided

\section*{Partitions}

\noindent\textbf{Instructions:} Not provided

\subsection*{Partition information}

\noindent\textbf{Partition name:} \\
JUPITER Booster Module

\vspace{0.5em}
\noindent\textbf{Requested amount of resources (node hours):} 25,000

\vspace{0.5em}
\noindent\textbf{Code(s) used:} Lux.jl, PyTorch, Neural Jump ODEs

\subsection*{Jobs}

\noindent\textbf{Number of jobs simultaneously:} 4 \\
\textbf{Wall clock time of a typical job execution (hours):} 24

\subsection*{Checkpoints}

\noindent\textbf{Are you able to write checkpoint?:} Yes \\
\textbf{Maximum time between 2 checkpoints (hours):} 4 \\
\textbf{Desirable maximum time between 2 checkpoints (hours):} 6

\subsection*{Cores/nodes}

\noindent\textbf{Average \# GPUs to be used per job:} 4 \\
\textbf{Maximum \# GPUs to be used per job:} 32 \\
\textbf{\# GPUs used per node:} 4 \\
\textbf{Maximum \# CPU cores per job:} 288 \\
\textbf{ID:} EHPC-AIF-2026SC01-062

\vspace{0.5em}
\noindent\textbf{3/7}

\vspace{0.5em}
\noindent\textbf{Maximum number of jobs running simultaneously on the GPU partition:} 8 \\
\textbf{Average \# CPU cores per job:} 128 \\
\textbf{\# of CPU cores used per node:} 72

\subsection*{Memory}

\noindent\textbf{Minimum job memory (total usage over all nodes in GB):} 100 \\
\textbf{Average job memory (total usage over all nodes in GB):} 300 \\
\textbf{Maximum job memory (total usage over all nodes in GB):} 2000

\subsection*{Storage}

\noindent\textbf{Maximum amount of SCRATCH needed at a time (TB):} 50 \\
\textbf{Maximum amount of WORK needed at a time (TB):} 10 \\
\textbf{Maximum amount of HOME needed at a time (TB):} 1 \\
\textbf{Maximum amount of ARCHIVE needed at a time (TB):} 50 \\
\textbf{Maximum \# files to be stored on SCRATCH (thousands):} 1000 \\
\textbf{Maximum \# files to be stored on WORK (thousands):} 100 \\
\textbf{Maximum \# files to be stored on HOME (thousands):} 10 \\
\textbf{Maximum \# files to be stored on ARCHIVE (thousands):} 500 \\
\textbf{Total amount of data to transfer to/from (TB):} 20

\vspace{0.5em}
\noindent\textbf{Justification of data transfer:} \\
Transfer of large retrospective 3D imaging datasets (PET/CT, MRI) from clinical PACS systems for training.

\vspace{0.5em}
\noindent\textbf{I/O Strategy:} \\
We utilize HDF5 with chunking for parallel I/O. Data is prefetched to GPU memory using `pin_memory=True` and multiple workers to saturate the NVLink-C2C bandwidth.

\subsection*{I/O}

\noindent\textbf{I/O data trafic R/W per hour:} 500 GB \\
\textbf{I/O files generated per hour:} 100 \\
\textbf{Specify if you would need to transfer the data to/from any of the following external data repositories or through data transfer federations:} None.

\section*{Code details, Development and Data management}

\noindent\textbf{Instructions:} Not provided

\vspace{0.5em}
\noindent\textbf{ID:} EHPC-AIF-2026SC01-062

\vspace{0.5em}
\noindent\textbf{4/7}

\subsection*{Workflows}

\noindent\textbf{Please describe the workflows you will be using:} \\
Our workflow involves: 1. Data Ingestion (DICOM to NIfTI conversion, anonymization). 2. Preprocessing (Spatial padding, normalization). 3. Model Training (Supervisors -> VAEs -> NJDEs). 4. Validation (Counterfactual generation). This pipeline is orchestrated via Slurm scripts.

\vspace{0.5em}
\noindent\textbf{Will you be using containerized solution?:} Yes (Apptainer).

\vspace{0.5em}
\noindent\textbf{If the project develops an AI model/method/software/workflow will it be available as open source?:} Yes.

\vspace{0.5em}
\noindent\textbf{Please provide a monthly plan of CPU/GPU resource usage by the project:} \\
Months 1-3: Data prep & Supervisor training (10% usage). Months 4-9: Heavy Generative/NJDE training (70% usage). Months 10-12: Validation & Refinement (20% usage).

\subsection*{Scalability and performance}

\noindent\textbf{Describe the scalability and performance of the application:} \\
We have benchmarked our framework using a "Super Heavy 3D ResNet-152" backbone (128x128x64 volume). We observed positive strong scaling from 1 GPU (40s/epoch) to 4 GPUs (12s/epoch), achieving a 3.3x speedup (82.5\% efficiency). Benchmarks were performed using the services and supercomputing resources of the \textbf{KI Servicezentrum Schleswig-Holstein \& KI (KISSKI)} (\url{https://kisski.gwdg.de/en/}).

\vspace{0.5em}
\noindent\textbf{Do you face bottlenecks in your AI solution? If yes select the type below::} Memory capacity (HBM).

\subsection*{Data details}

\noindent\textbf{Is the data used in the project open for communities?:} No (Patient Data).

\vspace{0.5em}
\noindent\textbf{Will the generated data by the project (if any) be open to other communities?:} Yes (Anonymized benchmarks).

\section*{Application Support Team (AST)}

\noindent\textbf{Instructions:} Not provided

\vspace{0.5em}
\noindent\textbf{Does your proposal require assistance from an AST on the selected partition(s)?:} No.

\section*{Collaboration and Funding}

\noindent\textbf{Instructions:} Not provided

\vspace{0.5em}
\noindent\textbf{Select one or more funding options applicable to this project:} European Commission (Horizon Europe).

\section*{Dissemination Strategy}

\noindent\textbf{Instructions:} Not provided

\vspace{0.5em}
\noindent\textbf{Dissemination strategy description:} \\
Publication in high-impact journals (Nature Medicine), open-source code release, and presentations at EANM/RSNA conferences.

\section*{Ethics Self-Assessment}

\noindent\textbf{Instructions:} Not provided

\subsection*{Respect for Human Agency}

\noindent\textbf{Please describe how your system ensures that end-users have the ability to control vital decisions about their own} \\
\textbf{ID:} EHPC-AIF-2026SC01-062

\vspace{0.5em}
\noindent\textbf{5/7}

\noindent\textbf{lives:} \\
The system is designed as a Decision Support System (DSS). The clinician always retains the final decision-making authority. The model provides explainable counterfactuals ("what-if" scenarios) to aid, not replace, human judgment.

\subsection*{Privacy \& Data Governance}

\noindent\textbf{Please describe how data is collected and processed from the aspect of lawfulness, fairness and transparency:} \\
Data is collected retrospectively under strict ethical approval. Patients are informed where possible, or waivers are obtained for anonymized retrospective studies.

\vspace{0.5em}
\noindent\textbf{What measures (such as anonymization, pseudonymisation, encryption, and aggregation) you took to safeguard the rights of data subjects?:} \\
All data is pseudonymized at source before upload to XNAT. DICOM tags containing PII are removed or hashed. Data is encrypted at rest and in transit.

\vspace{0.5em}
\noindent\textbf{Please describe the measures you employ to prevent data breaches and leakages:} \\
Access control via JuDoor/LDAP. Data is stored in secure, access-controlled directories on the supercomputer. No data leaves the secure environment.

\subsection*{Fairness}

\noindent\textbf{How do you ensure avoiding algorithmic bias, in input data, modelling and algorithm design?:} \\
We employ a multi-center dataset to ensure geographic diversity. We explicitly model and disentangle confounders (e.g., scanner type, age) using our Causal AI framework to prevent the model from learning spurious correlations.

\subsection*{Individual, and Social and Environmental Well-being}

\noindent\textbf{If relevant, describe how the AI system is mindful of all stakeholders and the environment:} \\
The project aims to reduce overtreatment (improving patient well-being) and optimize resource usage in hospitals. We use energy-efficient training strategies (Green AI) where possible.

\subsection*{Transparency}

\noindent\textbf{Are the end-users aware that they are interacting with an AI system?:} Yes.

\vspace{0.5em}
\noindent\textbf{Please describe how the participants and/or end-users will be informed about interacting with an AI system, and about its purpose, capabilities, limitations, benefits and risks:} \\
Clinicians using the system will undergo training on its capabilities and limitations (e.g., it is a support tool, not a diagnostic device). Output reports will clearly label AI-generated content.

\subsection*{Accountability}

\noindent\textbf{Please describe how your system ensures that potential ethically and socially undesirable effects will be detected, stopped, and prevented from reoccurring:} \\
We implement a "human-in-the-loop" validation step. Any erroneous or biased prediction flagged by clinicians is fed back into the system for retraining and bias correction.

\section*{Excluded Reviewers}

\subsection*{Excluded Reviewer \#1}

\noindent\textbf{Full Name:} Not provided \\
\textbf{E-mail:} Not provided \\
\textbf{Affiliation:} Not provided

\subsection*{Excluded Reviewer \#2}

\noindent\textbf{ID:} EHPC-AIF-2026SC01-062

\vspace{0.5em}
\noindent\textbf{6/7}

\noindent\textbf{Full Name:} Not provided \\
\textbf{E-mail:} Not provided \\
\textbf{Affiliation:} Not provided

\subsection*{Excluded Reviewer \#3}

\noindent\textbf{Full Name:} Not provided \\
\textbf{E-mail:} Not provided \\
\textbf{Affiliation:} Not provided

\section*{Data Consent}

\noindent\textbf{Instructions:} Not provided

\vspace{0.5em}
\noindent\textbf{space:} Not provided

\vspace{0.5em}
\noindent\textbf{In case the proposal is awarded, EuroHPC JU would like to publish the Principal Investigator’s and Team Members’ names and organizations. This may involve sharing this information on our website, social media channels, or in other promotional materials related to the project. Please provide your consent below.:} Yes, I consent.

\vspace{0.5em}
\noindent\textbf{In order to submit the proposal, you must accept the call's Terms of Reference available as a document at \url{https://www.eurohpc-ju.europa.eu/eurohpc-ju-call-proposals-ai-science-and-collaborative-eu-projects_en}. By ticking the box below you confirm that you have read and understood the terms and conditions of the call.:} Yes, I confirm.

\vspace{0.5em}
\noindent\textbf{space:} Not provided

\vspace{0.5em}
\noindent\textbf{In case the proposal is recommended to be awarded but due to the first-come-first-serve principle is not funded, do you agree on moving your proposal to the next cut-off period with a priority in funding?:} Yes.

\section*{Administrative self-assessment checklist}

\noindent\textbf{Administrative self-assessment checklist}

\vspace{0.5em}
\noindent\textbf{PI CV uploaded (most recent):} Yes.

\vspace{0.5em}
\noindent\textbf{Correct Project Scope and Plan template used:} Yes.

\vspace{0.5em}
\noindent\textbf{Maximum page limit of the Project Scope and Plan respected (10 pages):} Yes.

\vspace{0.5em}
\noindent\textbf{All sections of the Project Scope and Plan are completed:} Yes.

\vspace{0.5em}
\noindent\textbf{Total amount of resources in the milestones table (section 3.2) matches the resources requested in the online form (Partitions):} Yes.

\end{document}
